% Separate n=4 continuation. GPT-6 Astra Pro, /root/fixed_surface.
% 2026-11-22 UTC; effort not exposed. No changes to S01--S11.

\section{HD-F1-S12: the exact higher-dimensional two-form pairing}
\label{sec:f1-higher-lagrangian-obstruction}

Let $A$ be a polarized Weil abelian variety of dimension $2n$, $n\geq3$,
and let $Y\subset A$ be an actual smooth subvariety of dimension $n$
whose whole class is $[Y]=c\theta^n+w$, where $c>0$ and $\theta w=0$.
For each nonzero invariant holomorphic two-form $\Omega$ one has
\[
 \int_Y\Omega\wedge\overline\Omega\wedge\theta^{n-2}
 =c\int_A\Omega\wedge\overline\Omega\wedge\theta^{2n-2}>0.
 \tag{F1S12.1}
\]
Indeed the Gysin projection formula gives the equality, since
$\theta^{n-2}w=0$. In a $\theta$-unitary basis with
$\theta=i\sum_{j=1}^{2n}z_j\wedge\bar z_j$ and
$\Omega=\sum_{j<k}a_{jk}z_j\wedge z_k$, direct exterior multiplication gives
\[
 \Omega\wedge\bar\Omega\wedge\theta^{2n-2}
   =\frac{\sum_{j<k}|a_{jk}|^2}{(2n)(2n-1)}\theta^{2n}.
 \tag{F1S12.2}
\]
Thus the actual restriction $H^0(A,\Omega_A^2)\to H^0(Y,\Omega_Y^2)$
is injective. In particular such a $Y$ is not Lagrangian for an invariant
nondegenerate two-form. This computes a restriction map; it makes no
assertion about injectivity of the full Hochschild trace of a derived
object. For $n=2$ the factor $\theta^{n-2}$ is absent, so the vanishing
step above does not apply, exactly as in HD-F1-S10.

\section{HD-F1-S13: the actual smooth genus-five norm source}
\label{sec:f1-norm-fourfold-source}

There is a nonempty Zariski open subset $U\subset\mathcal M_5$ such that
for every connected etale cyclic triple cover $\pi:Z\to C$, $C\in U$,
the actual canonical norm fibre
\[
 X=\operatorname{Sym}^8Z,\quad Q=|K_C|=\mathbb P^4,\quad
 T=(\pi^{(8)})^{-1}Q=F^{-1}(K_C),\quad F=\operatorname{Nm}_\pi\circ a
 \tag{F1S13.1}
\]
is a disjoint union of three smooth connected fourfolds $T_j$. Each
maps finitely to $Q$ with degree $3^7$. Write
$\eta=[z_0+\operatorname{Sym}^7Z]$, $\vartheta=a^*\Theta_Z$,
$\nu=F^*\Theta_C$ and $p:T\to Q$. Then
\[
 \begin{gathered}
 0\longrightarrow p^*\mathcal O_Q(-1)\longrightarrow
 H^0(C,K_C)\otimes\mathcal O_T\longrightarrow N^*_{T/X}\longrightarrow0,\\
 N_{T/X}=p^*\Omega^1_{\mathbb P^4}(1),\qquad
 K_T=\vartheta+\eta\text{ numerically, and }K_T\text{ is ample}.
 \end{gathered}
 \tag{F1S13.2}
\]
All covers in the indicated actual open set are used; smoothness of
this norm fibre is proved below, rather than assumed.

\paragraph{1. Uniform smoothness over a nonempty open of genus-five curves.}
For a positive partition $\mathbf m=(m_1,\ldots,m_s)$ of eight, let
$\mathcal H_{\mathbf m}$ be the scheme of distinct marked points on a
smooth genus-five curve satisfying
$\mathcal O_C(\sum m_ix_i)=K_C$.
Polishchuk, \emph{Moduli spaces of curves with effective $r$-spin
structures}, Theorem 1.1(a), applied with $d=0$ (and any $r\geq2$),
proves that this stack is smooth of pure dimension $8+s$.
Its actual definition as the inverse image of the canonical section of
the universal Picard scheme is in the same paper, Section 2. Thus it
has exactly the scheme structure used here, not merely a parametrization
of the same closed points.

Work with a fine level cover to use schemes. There are finitely many
partitions and finitely many components in these finite-type strata.
Remove the closures of the images of the nondominating components.
Generic smoothness applied to every remaining component gives a common
nonempty open $U$ of $\mathcal M_5$ on which all these morphisms are
smooth. Remove the hyperelliptic locus as well.
For $C\in U$, every canonical divisor $E=\sum m_ix_i$ belongs to one
such fibre, of dimension $s-4$. Its tangent space is the kernel of the
weighted Abel differential
\[
 \bigoplus_i T_{x_i}C\longrightarrow H^1(C,\mathcal O_C).
\]
The dual is the weighted evaluation map of canonical forms at the
reduced divisor $E_{\rm red}$. All weights $m_i$ are nonzero in
characteristic zero. Hence its kernel has dimension
$s-5+h^0(K_C-E_{\rm red})$. Equality with the smooth fibre dimension
$s-4$ proves
\[
            h^0(C,K_C-E_{\rm red})=1
       \quad\text{for every }E\in|K_C|,\ C\in U.
 \tag{F1S13.3}
\]

For $D=\sum m_zz\in T$, put
$E_{\max}(D)=\sum_x\max_{z\mid x}m_z\,x$. The exact dual differential
of the norm Abel map, including all multiplicities, is
\[
 H^0(C,K_C)\xrightarrow{\pi^*}H^0(Z,K_Z)
                    \longrightarrow H^0(D,K_Z|_D),
 \quad\ker(dF_D)^*=H^0(C,K_C-E_{\max}(D)).
 \tag{F1S13.4}
\]
This follows in a local etale coordinate by requiring the appropriate
vanishing order at each of the three lifts. For $E=\pi_*D$ one has
$E_{\rm red}\leq E_{\max}\leq E$. Its canonical section belongs to
(F1S13.4), while (F1S13.3) bounds the dimension by one. Thus $dF_D$ has
rank four everywhere on $T$.
The map $\pi^{(8)}$ is finite flat of degree $3^8$: it is finite between
smooth varieties of the same dimension, so the local Cohen--Macaulay
flatness criterion applies. Its inverse image of $Q$ has pure dimension
four. The scheme fibre $F^{-1}(K_C)$ also has tangent dimension four at
every point, so it is regular and smooth. The kernel line in (F1S13.4)
is precisely the canonical section defining $p(D)$; its globalization
is the first map of (F1S13.2). The Euler sequence identifies its quotient
and the normal bundle as asserted.

\paragraph{2. Three actual components.}
For the canonical embedding $C\subset\mathbb P^4$, the incidence variety
$\{(z,H):\pi(z)\in H\}$ is a $\mathbb P^3$-bundle over connected $Z$.
Over hyperplanes meeting $C$ in eight distinct points its monodromy is
therefore transitive on the twenty-four lifted points. It is a subgroup
of $(\mathbb Z/3)^8\rtimes\mathfrak S_8$, commuting with the deck action.
A meridian around an ordinary tangent hyperplane is a colored
transposition: it exchanges two labels with colors $a,-a$ and fixes the
others. Such meridians and their conjugates generate the monodromy,
since projective four-space is simply connected and the nongeneric
part of the dual discriminant has complex codimension at least two
in the parameter space. Thus the monodromy lies in the sum-zero wreath
subgroup $G(3,3,8)$.

The graph of projected transpositions is connected. Gauge its colors
on a spanning tree to zero; its tree transpositions generate the plain
$\mathfrak S_8$. If every remaining color were zero, the twenty-four
labels would have three orbits, contrary to transitivity. A colored
transposition of nonzero color, multiplied by its plain transposition,
gives a nonzero vector $a(e_i-e_j)$ in the deck subgroup. Its
$\mathfrak S_8$-conjugates span the whole sum-zero subgroup. The actual
monodromy is consequently $G(3,3,8)$. On the $3^8$ choices of one lift
of each point it has three orbits, indexed by the sum of the eight
colors, each of size $3^7$. Since $T$ is smooth, its irreducible
components are disjoint and these are exactly its three connected
components. The diagonal deck generator adds eight to the sum modulo
three, so it cycles them. This is the same explicit monodromy argument
as HD-F1-T07, with eight labels in place of four.

\paragraph{3. Actual cycle class, normal class and exact intersections.}
On $\operatorname{Sym}^8C$ let $x,\theta$ be the point and theta classes.
The canonical plane $Q$ is the rank-four degeneracy locus of
$H^0(K_C)\otimes\mathcal O\to\Omega^1_{\operatorname{Sym}^8C}$.
The local Picard cohomology complex at $K_C$ is a row of five functions
whose derivative is the perfect Serre pairing. Analytic coordinates
make this row $(u_1,\ldots,u_5)$. In the projective-kernel chart $s_1=1$,
the cotangent minors generate $(u_2,u_3,u_4,u_5)$; thus the determinantal
scheme is the reduced $Q$ with its expected codimension four.
Porteous, together with the universal-divisor GRR calculation in
HD-F1-S09, gives
\[
 [Q]=x^4-x^3\theta+\frac{x^2\theta^2}{2}
                    -\frac{x\theta^3}{6}+\frac{\theta^4}{24},
\]
\[
 [T]=81\eta^4-27\eta^3\nu+\frac92\eta^2\nu^2
                         -\frac12\eta\nu^3+\frac1{24}\nu^4.
 \tag{F1S13.5}
\]
The exact endomorphism identity for the norm and theta polarizations
shows that $\nu$ has eigenvalues $3$ with multiplicity five and $0$ with
multiplicity eight, relative to $\Theta_Z$. The standard symplectic
exterior expansion, or pullback to $Z^8$ with its actual $8!$ quotient
degree, then gives
\[
 \int_X\eta^{8-a-b}\vartheta^a\nu^b
       =a!b!\binom{13-b}{a}\binom5b3^b\quad(a+b\leq8).
 \tag{F1S13.6}
\]
Combining (F1S13.5)--(F1S13.6) and the cyclic symmetry of the components
produces
\[
 \begin{array}{c|rrrrr}
  &\eta^4&\eta^3\vartheta&\eta^2\vartheta^2&\eta\vartheta^3&\vartheta^4\\\hline
 \displaystyle\int_{T_j}&27&216&1512&9072&45360.
 \end{array}
 \tag{F1S13.7}
\]
The hyperplane class $h=p^*\mathcal O_Q(1)$ equals $3\eta$ in cohomology.
For $g(Z)=13,d=8$, universal-divisor GRR gives
$\operatorname{ch}(T_X)=(-4-\vartheta)e^\eta+12$. Dividing the resulting
Chern class by $c(N_{T/X})=(1+3\eta)^{-1}$ yields
\[
             c(T_T)=(1+3\eta)(1+\eta)^{-4}
                            \exp\left(-\frac{\vartheta}{1+\eta}\right).
 \tag{F1S13.8}
\]
In particular $K_T=\eta+\vartheta$. The class $\eta$ is ample: under
the finite quotient $Z^8\to\operatorname{Sym}^8Z$ it pulls back to the
sum of the eight point-divisor classes, an ample external-product line
bundle. The class $\vartheta$ is nef, being pulled back from an ample
Jacobian theta divisor. Their sum is ample, proving the last assertion
of (F1S13.2).

For independent reproducibility the four Chern classes are
\[
\begin{aligned}
 c_1&=-\eta-\vartheta,\\
 c_2&=-2\eta^2+2\eta\vartheta+\vartheta^2/2,\\
 c_3&=10\eta^3-3\eta\vartheta^2/2-\vartheta^3/6,\\
 c_4&=-25\eta^4-10\eta^3\vartheta+3\eta^2\vartheta^2/2
                         +2\eta\vartheta^3/3+\vartheta^4/24.
\end{aligned}
 \tag{F1S13.9}
\]
Their component Chern numbers and Todd genus are
\[
\begin{array}{c|rrrrrrr}
 &c_1^4&c_1^2c_2&c_2^2&c_1c_3&c_4&\chi(\mathcal O)&c_4(-T_T)\\\hline
 \displaystyle\int_{T_j}&91611&53190&30888&22518&7371&318&594.
\end{array}
 \tag{F1S13.10}
\]
The last column is the virtual-normal number
$c_1^4-3c_1^2c_2+c_2^2+2c_1c_3-c_4$; it is not asserted to be the normal
bundle of an embedded image.

\section{HD-F1-S14: the whole character of the actual derived Abel image}
\label{sec:f1-norm-fourfold-character}

Translate a connected norm torsor to the Prym eightfold $P$, and let
$a_j:T_j\to P$ be the actual Abel morphism. Put
\[
                    \mathcal E_j=Ra_{j*}\mathcal O_{T_j}.
 \tag{F1S14.1}
\]
This is an actual bounded coherent perfect complex, since $a_j$ is
proper and $P$ is smooth. With $\beta=\Theta_Z|_P$, of polarization type
$(1,1,1,1,3,3,3,3)$, its complete rational character is
\[
 \operatorname{ch}(\mathcal E_j)=w_j+\frac{\beta^4}{72}
       -\frac{\beta^5}{120}+\frac{17\beta^6}{6480}
       -\frac{13\beta^7}{22680}+\frac{53\beta^8}{544320},
 \qquad 0\ne w_j\in W_K(P),\quad w_j^2=18.
 \tag{F1S14.2}
\]
Here $K=\mathbb Q(\sqrt{-3})$ acts by the actual deck automorphism.
In particular this is the required whole determinant-plus-polarization
character with $B=0$, without discarding the higher-degree terms.

\paragraph{1. Why the whole character has exactly these components.}
Looijenga, \emph{Prym representations of mapping class groups},
Geom. Dedicata 64 (1997), Proposition 2.2 and Theorem 2.4, computes the
image of the liftable mapping-class group on the primitive cyclic-cover
lattice as $U^\#(H_C)$. For our degree-three cover of a genus-five curve,
its connected Zariski closure is the split unitary group $SU(4,4)$:
the proposition identifies the cyclotomic lattice with four hyperbolic
pairs, and the elementary unitary transformations used in the proof of
the theorem give the corresponding unitary root subgroups in the
Zariski closure. Equivalently, the integral powers in each such root
group are Zariski dense in that root group, and these groups generate
the special unitary group. A finite level change or finite component
label cover does not change its connected Zariski closure. Removing a
proper algebraic closed locus from the base also does not change it,
since the inclusion of the open subset surjects on the topological
fundamental group.

The smooth proper family of the $T_j$ and its actual relative Abel map
make $\operatorname{ch}(Ra_{j*}\mathcal O_{T_j})$ a horizontal class.
This can be checked directly by relative GRR and proper cohomological
Gysin. A local translation trivializes each norm torsor, and translations
act trivially on cohomology, so there is no dependence on the chosen
trivialization. Thus every homogeneous component of this character is
invariant under the connected group just identified. Over $\mathbb C$,
$H^1(P)=V\oplus V^*$ for the standard eight-dimensional representation
of $SL_8$. In
$\bigwedge^{2k}(V\oplus V^*)=\bigoplus_{p+q=2k}\bigwedge^pV\otimes
\bigwedge^qV^*$, an invariant occurs for $p=q$, giving $\beta^k$,
and additionally for $(p,q)=(8,0),(0,8)$, giving the two determinant
lines at $k=4$. This follows from irreducibility and inequivalence of
the exterior powers, with $\bigwedge^0V=\bigwedge^8V$ the only duplicate.
The support of the derived pushforward has dimension at most four,
so $\operatorname{ch}_k=0$ for $k<4$. We have proved the full shape in
(F1S14.2) by monodromy of this actual twelve-dimensional cover family;
it does not require this family to dominate the sixteen-dimensional
Weil domain.

GRR has no denominator from the ambient tangent bundle, which is
trivial. Hence
\[
 \int_P\operatorname{ch}_{4+j}(\mathcal E_j)\beta^{4-j}
       =\int_{T_j}\operatorname{td}_j(T_{T_j})\vartheta^{4-j}.
 \tag{F1S14.3}
\]
For $j=0,1,2,3,4$ the right sides are respectively
$45360,-27216,8568,-1872,318$. Divide by
$\beta^8=8!\,3^4=3265920$ to obtain exactly the five coefficients in
(F1S14.2). The determinant terms pair trivially with $\beta$ by their
exterior description. The exact arithmetic certificate
\texttt{check\_norm\_fourfold.py} checks all these identities.

\paragraph{2. The exact determinant norm from Schoen's actual quotient.}
Let $\sigma$ be the deck generator and set
\[
 Q_0=3\operatorname{Id}-(1+\sigma+\sigma^2),\qquad
 f:Z^8\xrightarrow{\sum a_Z}JZ\xrightarrow{Q_0}P.
 \tag{F1S14.4}
\]
The image curve generates $P$, so its eighth sum dominates $P$ and $f$
is generically finite. Let
$q:Z^8\to W_0=Z^8/((\mathbb Z/3)^7\rtimes\mathfrak S_8)$.
This is Schoen's actual sum-zero quotient, of degree
$D_0=3^7\,8!$. Its induced map to $\operatorname{Sym}^8C$ is etale of
degree three. The inverse image of $Q$ consists of disjoint copies
$Q_0',Q_1',Q_2'$ of $\mathbb P^4$. Each has normal bundle
$\Omega^1_{\mathbb P^4}(1)$ and self-intersection $+1$, by the same
Euler-sequence calculation (or Schoen 1988, Lemma 2.6, unramified case).
For a primitive cube root $\zeta$ put
$z_\chi=\sum_j\zeta^{-j}[Q_j']$ and $u_\chi=q^*z_\chi$.
Then $z_\chi z_{\chi^{-1}}=3$ and
\[
                  \int_{Z^8}u_\chi u_{\chi^{-1}}=3D_0.
 \tag{F1S14.5}
\]
Here is a direct identification of these character classes with
$f^*W_K(P)$, also matching Schoen's unramified calculation. The class
$u_\chi=q^*z_\chi$ is invariant under the sum-zero deck subgroup and
has a nontrivial character under its cyclic quotient. In the Kunneth
expansion of $H^8(Z^8)$, these conditions force the character on each
of the eight factors to be the same nontrivial character. Each factor
therefore lies in $H^1(Z)_\chi$, since $H^0$ and $H^2$ have trivial
deck character. Permutation invariance, with the Koszul sign for
permuting degree-one classes, identifies the invariant subspace with
$\bigwedge^8H^1(Z)_\chi$, a line. Pull back the eight forms in this
eigenspace by $f$ and expand over the factors. Two same-character
forms on one curve have zero product, so the surviving terms are
exactly the alternating tensors spanning this line. The operator
$Q_0$ acts as three on these eigenspaces. Hence the actual class
$u_\chi$ belongs to $f^*W_K(P)$, with no additional normalization
premise. The other character is identical. This is the two-line
primitive-character subspace in Schoen's Sections 1--2.

The theta-norm identities give
$Q_0^*\beta=9\Theta_Z-3\operatorname{Nm}^*\Theta_C$ on $JZ$.
Its relative eigenvalues are zero on the five-dimensional invariant
part and nine on the eight-dimensional Prym part. Pullback by the
ordered sum consequently gives
\[
 \deg f=\frac{(8!)^2 9^8}{8!3^4}=8!3^{12}.
 \tag{F1S14.6}
\]
The class $a_{j*}[T_j]$ has positive degree $45360$, hence its image
$Y_j$ has dimension four. Every nonempty fibre of the Abel morphism is
a whole complete linear system $|L|$, because $\operatorname{Nm}L=K_C$
already implies that all its effective divisors lie in $T$.
A finite such fibre is a single point, so $a_j$ has generic degree one
and $a_{j*}[T_j]=[Y_j]$. On each connected norm torsor, $Q_0$ is
multiplication by three followed by translation. The actual ordered
quotient has degree $8!$, and $[3]_*$ on middle degree-eight classes of
an eightfold is multiplication by $3^8$. Therefore
\[
               f_*u_\chi=8!3^8\sum_j\zeta^{-j}[Y_j].
 \tag{F1S14.7}
\]
The left side belongs to $W_K(P)$, and the Lefschetz terms cancel on the
right. The deck generator permutes $w_j$ and acts on the two determinant
lines with eigenvalues $\zeta,\zeta^2$ (their characters are the eighth
powers of the nontrivial characters). Consequently the pairing of the
two character sums is $(9/2)w_j^2$. Since $u_\chi\in f^*W_K(P)$ and
$f_*f^*=(\deg f)\operatorname{Id}$, (F1S14.5)--(F1S14.7) give
\[
 (8!3^8)^2\frac92w_j^2=(8!3^{12})(3\cdot3^7 8!),
          \qquad w_j^2=18.
 \tag{F1S14.8}
\]
In particular the component is nonzero and
\[
                 [Y_j]^2=\frac{\beta^8}{72^2}+18=648.
 \tag{F1S14.9}
\]

\section{HD-F1-S15: the actual Abel image is singular}
\label{sec:f1-norm-fourfold-singular-image}

For every cover in the open set of HD-F1-S13, each actual reduced Abel
image $Y_j\subset P$ is singular, and the morphism $a_j:T_j\to Y_j$
has positive-dimensional complete-linear-system fibres.

First, $a_j$ is proper and birational by the exact fibre argument above.
If it were an embedding, the tangent-normal exact sequence in the
abelian eightfold would imply
\[
 [Y_j]^2=\int_{T_j}c_4(N_{T_j/P})
               =\int_{T_j}c_4(-T_{T_j})=594,
\]
contradicting (F1S14.9). Even without the exact norm calculation, the
primitive $(4,4)$ polarization pairing gives $w_j^2\geq0$, so
$[Y_j]^2\geq630>594$ is already a contradiction.
At any divisor $D$ with $h^0(\mathcal O_Z(D))=1$, the divisor exact
sequence shows that the Abel differential is injective: its kernel is
$H^0(\mathcal O_Z(D))/H^0(\mathcal O_Z)$. If all fibres were finite,
all these kernels would vanish; a proper injective immersion would give
an embedding. Thus at least one complete system $|L|$ is positive
dimensional on each component.

Suppose its image $Y_j$ were smooth. The birational morphism between
smooth varieties has the effective exceptional Jacobian divisor
$R=K_{T_j}-a_j^*K_{Y_j}$. Since $K_{T_j}$ is ample, $R$ is $a_j$-nef
(in fact it has positive degree on each contracted curve).
The negativity lemma for an effective exceptional divisor on a
projective birational morphism forces $R=0$.
For an exact primary locator see Kollar--Mori, \emph{Birational geometry
of algebraic varieties}, Lemma 3.39(1): apply it to $-R$, whose negative
is relatively nef and whose direct image is zero. The same application
is recorded in Greb--Kebekus--Peternell, \emph{Projectively flat klt
varieties}, Section 2.10, proof of Lemma 2.11.
With $R=0$, the determinant of the differential nowhere vanishes, so
$a_j$ is etale; a proper birational etale morphism is an isomorphism.
This contradicts the preceding nonembedding calculation. The actual
Abel image is therefore singular.

After shrinking to a further nonempty open of the cover moduli, all its
positive-dimensional Abel fibres are $\mathbb P^1$. Indeed Schwarz,
\emph{Brill--Noether theory for cyclic covers}, Theorem 1.1, proves that
a general etale cyclic cover has no $g_d^r$ when
$\rho(\widetilde g,r,d)<-r$. Here
$\rho(13,2,8)=13-3(13-8+2)=-8<-2$, so $h^0(L)\leq2$ for all degree-eight
line bundles. This is an actual consequence of that theorem, not an
assumption about the singularities. No finiteness or normal-bundle
splitting of these pencil fibres is asserted here.

\paragraph{The exact map remaining for the derived trace calculation.}
On the open subset of $T_j$ where $h^0(L)=1$, let
$\mathcal V=\ker(H^0(Z,K_Z)\otimes\mathcal O\to K_Z^{[8]})$.
It is a rank-five vector bundle. Restriction of the ambient cotangent
space followed by the actual differential gives
\[
 0\longrightarrow p^*\mathcal O_{\mathbb P^4}(-1)
 \longrightarrow\mathcal V\longrightarrow N^*_{T_j/P}\longrightarrow0.
 \tag{F1S15.1}
\]
The first line is the canonical norm section from (F1S13.4).
To see the quotient directly, the conormal of the Abel immersion in
$JZ$ is $H^0(K_Z-D)$, while the cotangent space of the norm torsor is
$H^0(K_Z)/\pi^*H^0(K_C)$. Their intersection is exactly the displayed
canonical line. Thus (F1S15.1) is the actual conormal of the immersive
open locus. It does not extend by assertion across the exceptional
Abel fibres or the singular image.

The specified perfect complex in (F1S14.1) has the whole character
(F1S14.2). The calculation above has not computed
$\operatorname{Ext}^{<0}(\mathcal E_j,\mathcal E_j)$, nor the kernel of
its actual Hochschild semiregularity trace
$\operatorname{Ext}^2(\mathcal E_j,\mathcal E_j)\to HH_{-2}(P)$,
with or without a specified finite-group invariant projector.
These are concrete remaining gates for this very object. In particular
no inference from the singularity of $Y_j$ to failure of that trace is
made. The exact smooth embedded argument of S10 is blocked both by
(F1S12.1) and by the actual singularity just proved; the derived object
and its trace retain their separate explicitly stated maps.
